\section{Kết luận}

Đồ án đã nghiên cứu bài toán phát hiện, theo dõi và đếm phương tiện từ dữ liệu
video giao thông, trong đó trọng tâm là xây dựng một chuỗi xử lý dựa trên học
sâu để xác định vị trí phương tiện trong từng khung hình và duy trì định danh
của chúng qua thời gian. Phương pháp được đề xuất sử dụng YOLOv8n làm bộ phát
hiện phương tiện và kết hợp với thuật toán theo dõi BYTE để liên kết các phát
hiện giữa các khung hình. Cách tiếp cận này tận dụng ưu điểm của YOLOv8n về tốc
độ và kích thước mô hình nhỏ, đồng thời khai thác khả năng của BYTE trong việc
sử dụng cả các phát hiện có độ tin cậy thấp nhằm hạn chế mất dấu đối tượng.

Về mặt thực nghiệm, mô hình YOLOv8n sau khi được fine-tune trên tập dữ liệu
UA-DETRAC cho thấy khả năng hội tụ ổn định. Các hàm mất mát trên tập huấn
luyện và tập xác thực đều giảm dần, trong khi các chỉ số Precision, Recall và
mAP có xu hướng cải thiện trong quá trình huấn luyện. Kết quả đánh giá phát
hiện cho thấy mô hình hoạt động tốt đối với các lớp có nhiều mẫu như xe hơi và
xe buýt, nhưng còn hạn chế đối với các lớp có ít dữ liệu hơn như xe tải nhỏ và
lớp khác. Điều này chứng minh rằng mô hình đã học được các đặc trưng hữu ích
của phương tiện trong dữ liệu giao thông, đồng thời cũng cho thấy sự mất cân
bằng dữ liệu vẫn là một yếu tố ảnh hưởng đáng kể đến hiệu năng phát hiện.

Đối với bài toán theo dõi, đồ án đã so sánh hai thuật toán SORT và BYTE trên
tập xác thực. Kết quả cho thấy BYTE đạt hiệu năng cao hơn SORT trên tất cả các
độ đo chính, bao gồm \HOTA, \MOTA\ và \IDFone. Đặc biệt, BYTE đạt \AssA\ và
\IDFone\ cao hơn, cho thấy khả năng duy trì danh tính đối tượng tốt hơn trong
quá trình theo dõi. Trên tập kiểm tra, hệ thống YOLOv8n kết hợp BYTE đạt
\HOTA\ bằng $60.716$, \DetA\ bằng $59.184$, \AssA\ bằng $62.800$, \MOTA\ bằng
$64.676$ và \IDFone\ bằng $70.345$. Các kết quả này chứng minh rằng việc kết
hợp bộ phát hiện YOLOv8n với BYTE là một hướng tiếp cận khả thi cho bài toán
theo dõi phương tiện trong video giao thông.

Các phân tích theo ngưỡng \HOTA\ cũng cho thấy hiệu năng tổng thể của hệ thống
chịu ảnh hưởng lớn từ chất lượng phát hiện. Khi ngưỡng đánh giá trở nên nghiêm
ngặt hơn, \DetA\ giảm nhanh hơn \AssA, cho thấy sai lệch về vị trí hoặc kích
thước hộp giới hạn là nguyên nhân quan trọng làm giảm chất lượng theo dõi. Như
vậy, đồ án không chỉ chứng minh được lợi ích của BYTE so với SORT trong việc
duy trì định danh, mà còn chỉ ra rằng cải thiện bộ phát hiện là hướng quan
trọng để nâng cao hiệu năng của toàn bộ hệ thống.

Ưu điểm chính của phương pháp đề xuất là cấu trúc đơn giản, dễ triển khai và có
khả năng cân bằng giữa tốc độ xử lý và độ chính xác. YOLOv8n có kích thước nhỏ,
phù hợp với các hệ thống có tài nguyên tính toán hạn chế, trong khi BYTE giúp
tận dụng tốt hơn các phát hiện ở nhiều mức độ tin cậy khác nhau. Bên cạnh đó,
việc tách riêng bộ phát hiện và bộ theo dõi giúp hệ thống linh hoạt, có thể thay
thế hoặc cải tiến từng thành phần mà không cần thay đổi toàn bộ kiến trúc.

Tuy nhiên, phương pháp vẫn còn một số hạn chế. Hiệu năng phát hiện chưa đồng
đều giữa các lớp do dữ liệu mất cân bằng, đặc biệt đối với các lớp ít mẫu như
\textit{others}. Các đối tượng nhỏ, bị che khuất hoặc xuất hiện ở điều kiện
quan sát khó vẫn có thể bị bỏ sót hoặc định vị chưa chính xác, từ đó ảnh hưởng
đến kết quả theo dõi. Ngoài ra, hệ thống hiện chủ yếu được đánh giá trên bài
toán phát hiện và theo dõi; phần đếm phương tiện vẫn phụ thuộc vào chất lượng
quỹ đạo đầu ra và cần được đánh giá sâu hơn trong các kịch bản giao thông thực
tế.

Trong tương lai, có thể cải thiện hệ thống theo một số hướng. Thứ nhất, cần xử
lý tốt hơn sự mất cân bằng dữ liệu bằng các chiến lược tăng cường dữ liệu, lấy
mẫu lại hoặc điều chỉnh hàm mất mát cho các lớp ít mẫu. Thứ hai, có thể thử
nghiệm các biến thể YOLO lớn hơn hoặc các kỹ thuật tối ưu suy luận để nâng cao
độ chính xác trong khi vẫn đảm bảo tốc độ xử lý. Thứ ba, việc bổ sung các đặc
trưng ngoại quan hoặc các thuật toán theo dõi mạnh hơn có thể giúp giảm nhầm
lẫn định danh trong các tình huống che khuất. Cuối cùng, cần đánh giá trực tiếp
độ chính xác của bài toán đếm phương tiện trên các vùng hoặc đường đếm cụ thể
để hoàn thiện hơn hệ thống phục vụ giám sát giao thông tự động.
